\section{Краевые задачи для дифференциальных уравнений в частных производных}
\label{sec:boundary_value_problems_for_pdes}

\subsection{Классификация дифференциальных уравнений в частных производных}

Общий вид линейного дифференциального уравнения с частными производными второго порядка относительно неизвестной функции \( u(x, y) \):

$$A \frac{\partial^2 u}{\partial x^2} + 2B \frac{\partial^2 u}{\partial x \partial y} + C \frac{\partial^2 u}{\partial y^2} + D \frac{\partial u}{\partial x} + E \frac{\partial u}{\partial y} + Gu = F.$$

Коэффициенты \( A, B, C, D, E, G \) и правая часть \( F \) есть некоторые функции независимых переменных \( x \) и \( y \). Если \( F \equiv 0 \), то уравнение называется однородным; в противном случае — неоднородным.

\begin{itemize}
	\item $B^2 - AC < 0$ - уравнение эллиптического типа (уравнения Лапласа, Пуассона, Гельмгольца);
	\item $B^2 - AC > 0$ - уравнение гиперболического типа (волновое и телеграфное уравнения);
	\item $B^2 - AC = 0$ - уравнение параболического типа (уравнение теплопроводности);
	\item если $B^2 - AC$ не имеет постоянного знака - уравнение смешанного типа.
\end{itemize}

\subsection{Начальные и граничные условия}

\textbf{Начальное условие нулевого порядка}
$$u(x, 0) = \varphi(x).$$

\textbf{Начальное условие первого порядка}
$$u_t(x, 0) = \psi(x).$$

\dotfill

\textbf{Условие Дирихле} — задаётся значение функции:

$$u|_{\partial\Omega} = g.$$

\textbf{Условие Неймана} — задаётся нормальная производная:

$$\frac{\partial u}{\partial n}|_{\partial\Omega} = g.$$

\textbf{Условие Робина} — линейная комбинация функции и её производной:

$$\alpha u + \beta \frac{\partial u}{\partial n} = g.$$

Также могут использоваться смешанные условия, когда на разных частях границы задаются разные типы условий.

\subsection{Эллиптические краевые задачи}

\textbf{Уравнение Лапласа (задача Дирихле)}

Найти функцию \( u(x,y) \), удовлетворяющую:
\[
\Delta u = 0 \quad \text{в} \quad \Omega,
\]
\[
u|_{\partial\Omega} = g(x,y).
\]

\textbf{Уравнение Пуассона (задача Неймана)}
\[
-\Delta u = f \quad \text{в} \quad \Omega,
\]
\[
\frac{\partial u}{\partial n} = g \quad \text{на} \quad \partial\Omega,
\]
при условии согласования
\[
\int_\Omega f \, d\Omega = \int_{\partial\Omega} g \, ds.
\]

\subsection{Параболические начально-краевые задачи}

\textbf{Уравнение теплопроводности (одномерное)}
\[
\frac{\partial u}{\partial t} = k \frac{\partial^2 u}{\partial x^2}, \quad x \in (0, L), \; t > 0,
\]
с начальным условием:
\[
u(x, 0) = \varphi(x),
\]
и граничными условиями (например, Дирихле):
\[
u(0, t) = u(L, t) = 0.
\]

\textbf{Уравнение конвекции–диффузии (параболическое)}
\[
\frac{\partial u}{\partial t} + \mathbf{v} \cdot \nabla u = D \Delta u, \quad x \in \Omega, \; t > 0,
\]
\[
u(x, 0) = \varphi(x),
\]
\[
\frac{\partial u}{\partial n} = 0 \quad \text{на} \quad \partial\Omega.
\]

\subsection{Гиперболические начально-краевые задачи}

\textbf{Волновое уравнение (одномерное)}
\[
\frac{\partial^2 u}{\partial t^2} = c^2 \frac{\partial^2 u}{\partial x^2}, \quad x \in (0, L), \; t > 0,
\]
с начальными условиями:
\[
u(x, 0) = \varphi(x), \quad \frac{\partial u}{\partial t}(x, 0) = \psi(x),
\]
и граничными условиями (например, Дирихле):
\[
u(0, t) = u(L, t) = 0.
\]

\subsection{Корректность постановки задачи}

В достаточном общем виде для операторного уравнения  
$$ A y = f, $$ 
где \( A : Y \to F \) — некоторый оператор, действующий из метрического пространства \( Y \) в метрическое пространство \( F \) (обычно \( Y \) и \( F \) — полные метрические пространства, в частности, банаховы), корректность вводится следующим образом.  

\textbf{Определение.} 
Задача нахождения элемента \( y \in Y \) по заданному элементу \( f \in F \) из уравнения называется \textbf{корректно поставленной по Адамару} для тройки \( (Y, A, F) \) или просто \textbf{корректной}, если:
\begin{enumerate}
	\item при каждом \( f \in F \) существует решение \( y \in Y \);
	\item это решение \( y \) единственно в \( Y \);
	\item решение непрерывно зависит от правой части уравнения (т.е. из сходимости \( f_n \to f \) по метрике пространства \( F \) следует сходимость \( y_n \to y \) по метрике пространства \( Y \)). 
\end{enumerate}

При нарушении любого из этих трех условий задача считается \textbf{некорректной}.

\subsection{Основные численные методы решения}
\begin{itemize}
	\item Метод конечных разностей;
	\item Метод конечных элементов;
	\item Метод конечных объемов;
	\item Метод прямых;
	\item Метод характеристик.
\end{itemize}

\subsection{Что важно сказать на экзамене}

Привести общий вид уравнения, дать классификацию уравнений, привести начальные и граничные условия. Понимать корректность задачи, знать пару постановок задач.

\newpage